\section{Constructed stateful training: protocol and accounting}
\label{app:stateful-training}
\label{app:stateful}

\paragraph{Environment.}
There are $n=H/4$ components in a directed dependency chain. Initially their
source, cached artifacts, attestations, and publications all reflect version
zero. The task changes the required version of one component, sampled
uniformly from the first $\lceil n/2\rceil$ positions, to one. Each component is
visited once in each of four rounds, in topological order. Every visit permits
one of five actions: no operation, edit source to the specification, build,
attest, or publish. These are actual tool-state transitions, including legal
no-ops. A build stores the current built parent prefix followed by the local
source version; attestations and publications retain immutable fingerprints.
Consequently, updating an upstream component does not automatically refresh
its dependants.

The reference verifier executes four constraints per component: the source
matches its requested version, and the artifact, attestation, and publication
each match the required complete prefix. The cheap verifier checks only the
local version of each published artifact. Its failures are conclusive
reference failures; every baseline uses this common shortcut. The observation
contains 15 features: an intercept, four phase indicators, local source and
requested versions, five inspectable metadata comparisons, and normalized
position, remaining horizon, and problem size. Constructing this observation
does not invoke the reference verifier. Public task requirements are not
hidden supervision; the purchased terminal labels are isolated from action
selection and unqueried estimator inputs.

\paragraph{Policy and optimization.}
A shared linear softmax maps these features to five action probabilities.
All methods initialize with probability $0.9$ on the round-appropriate action
and $0.025$ on each other action. This scripted prior is disclosed because a
nearly uniform policy would yield little terminal success at larger $H$.
The initial greedy policy already follows a correct workflow. We execute
300 updates per run, with eight newly sampled trajectories per update and
a learning rate of $0.5$. For trajectory $i$, the score is the mean action
score, $H^{-1}\sum_{t=1}^{H}\nabla_\theta\log\pi_\theta(a_{it}\mid o_{it})$.
The estimated group coefficients weight these scores in one on-policy SGD
update. There are no Adam updates, replay epochs, learned critics, or auxiliary
rewards. The norm cap is one; it activates in none of the 54,000 updates
(maximum observed norm $0.642$). Positive and negative coefficients are added
at likelihood ratio one, so these runs do not test off-policy clipping.

\paragraph{Auditing rules.}
After cheap-negative certification, $M$ candidate labels remain in their
original order. The completion model assigns them probability $0.55$ and
assigns certified negatives probability zero. It is fixed and not fitted on
training or reporting labels. Whole-group residual correction buys all $M$
labels with probability $1/2$ and otherwise uses the same completion vector.
Sequential completion (\texttt{rr\_completion} in the released code) uses
\[
 Q_j=\operatorname{clip}\bigl(c j^{-1/2},0.005,1\bigr),
 \qquad \sum_{j=1}^{M}Q_j=M/2,
\]
with the scale $c$ chosen before observing any purchased label. The empty
candidate set needs no purchases. Reward HT independently buys candidates
with probability $1/2$, substitutes twice the observed reward and zero for
unbought candidates, and then normalizes. Partial replacement uses the same
inclusion probabilities but retains unbought cheap labels. Both are distinct
from correction of the final group coefficients.

A subsequent outcome-free diagnostic evaluated the model-based increment
second moments of the homogeneous completion prior. For every
$M\in\{1,\ldots,8\}$, the monotone schedule minimizing the model's Euclidean
coefficient error was $Q_j=1/2$ for every candidate. It therefore coincided
with whole-group residual correction. This is a boundary case of the
allocation problem, not a reason to replace the frozen primary schedule or
a claim about the true policy-gradient error.

\paragraph{Randomness and evaluation.}
Ten independent training seeds are used for each horizon and method. Within
a seed, task and action uniforms are paired across methods; learned policies
can still produce different subsequent trajectories. Audit randomness uses a
separate stream. Reporting uses 256 stochastic episodes at update zero and
every 25 updates, with a separate stream reused for pairing across checkpoints
and methods. Reporting labels do not select parameters, schedules, updates, or
stopping rules. These are independent episode draws from the declared task
family, not unseen repository tasks. The horizon and chain size co-vary, so
their individual effects are not identified.

For SC minus whole-group residual, the endpoint differences and pointwise
95\% Student-$t$ intervals over ten paired training seeds are, in percentage
points, $0.078\;[-0.578,0.734]$ at $H=12$,
$-0.195\;[-1.110,0.720]$ at $H=24$, and
$1.328\;[-0.297,2.953]$ at $H=48$. These intervals use nine degrees of freedom,
are not adjusted for multiple comparisons, and are not non-inferiority tests.

\paragraph{Cost and saved evidence.}
The complete experiment executes 432,000 training trajectories and 12,096,000
tool actions in 259.07 seconds on the local CPU. Python 3.12.14 and NumPy 2.3.5
are recorded in the run manifest. Training purchases 185,536 reference checks
across all methods. Separately, reporting performs 599,040 checks, and
post-training replay performs 432,000 diagnostic checks. These diagnostic
labels are computed only after each entire optimization run; they never
enter its updates. One reference query executes $4n$ constraints. Query counts
are resource units; no artificial latency is added. The measured training
times in the main text exclude reporting and post-training diagnostics and
are sums over the 30 runs per method. They are shared-host diagnostics with a
fixed method order, not an isolated randomized systems benchmark.

The primary comparison fixes rollout count. A secondary released table selects
recorded checkpoints below a common descriptive query ceiling and reports
actual query and rollout counts. It does not interpolate, leaves unequal
unspent budget, and is not an exactly matched-cost experiment. We do not use
it to claim equal overall compute.

Every training action sequence, task change location, cheap label, purchased
label and mask, coefficient vector, and policy checkpoint is retained in
compressed arrays. Immutable tool replay reconstructs terminal states and
checks purchased labels against post-training truth. All 180 trajectory banks,
five numerical CSV files, and four executed-source snapshots pass the recorded
ledger and hash checks. Independent CPU checks compare trajectory scores with
held-action finite differences (maximum error $1.40\times10^{-11}$), and
integrate every sequential stopping cutoff on generated release groups
(maximum coefficient error $4.44\times10^{-16}$). These checks do not claim a
second complete execution of all 180 training runs.
